(1) 两场相加后强度交叉项为 \(\cos[(\mathbf k_1-\mathbf k_2)\cdot\mathbf r+\phi_0]\)，故 \(\mathbf q=\mathbf k_1-\mathbf k_2\)，其模长 \(|\mathbf q|=2k\sin\theta\)。相邻明纹间距 \(d=2\pi/|\mathbf q|=\lambda/(2\sin\theta)\)，条纹法向沿 \(\mathbf q\)。
(2) 轨迹 \(\mathbf r(t)=\mathbf r_0+\mathbf vt\) 代入交叉相位，得到角拍频 \(\Omega=\mathbf q\cdot\mathbf v\)，普通频率 \(f=\Omega/(2\pi)=v_\perp/d\)。符号取决于穿越方向；只测强度功率谱时正负号可能无法区分。
(3) 线性映射 \(f=v_\perp/d\) 将Gaussian速度分布映为Gaussian频率分布，中心 \(\bar f=\bar v/d\)、标准差 \(\sigma_f=\sigma_v/d\)，频率半高全宽 \(2\sqrt{2\ln2}\,\sigma_f\)。有限照明时间还产生额外傅里叶展宽，不能误认为纯速度分布。
(4) 只有沿 \(\mathbf q\) 的速度分量参与调制，垂直分量不可测。\(\theta\to0\) 时 \(d\to\infty\)、\(f\to0\)，测速灵敏度消失；两束互不相干时交叉项平均为零，条纹和可用拍频一同消失。
